\section{Basic Concepts of Rings}

Good examples of rings to keep in mind are \(\ZZ, \ZZ / n, \RR \brac{x}\).

\begin{definition}[Ring]
    \label{def:ring}%
    A \emph{ring} is a quintuple \(\para{R, +, \cdot, 0_R, 1_R}\) where \(R\) is a set, \(0_R, 1_R \in R\), and
    \begin{align*}
        \functiondomain{+&}{R \times R}{R} && \text{``addition''} \\
        \functiondomain{\cdot&}{R \times R}{R}, && \text{``multiplication''}
    \end{align*}
    such that
    \begin{itemize}
        \item \(\para{R, +, 0_R}\) is an abelian group;
        \item the operation \(\cdot\) is associative: for all \(a, b, c \in R\), \(a \cdot \para{b \cdot c} = \para{a \cdot b} \cdot c\);
        \item \(1_R\) is a multiplicative identity for \(R\): for all \(a \in R\), \(1_R \cdot a = a \cdot 1_R = a\);
        \item multiplication is distributive over addition: for all \(r_1, r_2, r_3 \in R\), \(r_1 \cdot \para{r_2 + r_3} = r_1 \cdot r_2 + r_1 \cdot r_3\), and \(\para{r_1 + r_2} \cdot r_3 = r_1 \cdot r_3 + r_2 \cdot r_3\).
    \end{itemize}
\end{definition}

\begin{notation}
    We will denote this as ``\(R\) is a ring''. The symbol \(-r\) will mean the additive inverse of \(r\), and note \(r - s = r + \para{-s}\).
\end{notation}

\begin{remark}
    In this course, all rings are assumed to \emph{commutative}, that \(r_1 \cdot r_2 = r_2 \cdot r_1\) for all \(r_1, r_2 \in R\).
\end{remark}

\begin{definition}[Subring]
    \label{def:subring}%
    A \emph{subring} of a ring \(R\) is a subset \(S \subseteq R\), such that \(0_R, 1_R \in S\), and if \(s, s' \in S\), then \(s + s', s \cdot s'\) also lie in \(S\), and \(\para{S, +, \cdot, 0_R, 1_R}\) is a ring.

    We note this as \(S \subring R\).
\end{definition}

\begin{example}[Examples of Rings]
    \label{ex:rings}%
    We have (the familiar number sets)
    \[
        \ZZ \subring \QQ \subring \RR \subring \CC
    \]
    are all rings.

    The \emph{Gaussian Integers} is
    \[
        \ZZ \brac{i} = \set{a + ib}{a, b \in \ZZ} \subring \CC.
    \]

    Similarly,
    \[
        \QQ \brac{\sqrt{2}} = \set{a + b \sqrt{2}}{a, b \in \QQ} \subring \RR.
    \]

    Less familiarly,
    \[
        \CC \ang{\epsilon} = \set{a + b\epsilon}{a, b \in \CC}
    \]
    with operations
    \[
        \para{a + b\epsilon} + \para{c + d\epsilon} = \para{a + c} + \para{b + d}\epsilon
    \]
    and
    \[
        \para{a + b\epsilon} \cdot \para{c + d\epsilon} = ac + \para{ad + bc} \epsilon,
    \]
    with
    \[
        0 = 0 + 0 \cdot \epsilon, 1 = 1 + 0 \cdot \epsilon.
    \]

    This is the ring of \emph{dual numbers}. In particular, \(\epsilon^2 = 0\) even though \(\epsilon \neq 0\).
\end{example}

\begin{definition}[Unit]
    \label{def:unit}%
    An element \(u \in R\) is a \emph{unit} if there exists \(v \in R\), such that \(uv = 1\).
\end{definition}

\begin{remark}
    Note that \(2 \in \ZZ \subring \QQ\), and \(2\) is a unit in \(\QQ\), but \textbf{not in \(\ZZ\)}. A unit in a ring is not necessarily a unit in a subring: \textbf{subrings do not preserve units}.
\end{remark}

\begin{definition}[Field]
    \label{def:field}%
    A ring \(R\) is a \emph{field} if every non-zero element is a unit.
\end{definition}

Note that the product of units must be a unit, although their sum might not be. The units form a subgroup under multiplication.